\documentclass[UTF8,11pt,a4paper]{ctexart}
\usepackage[a4paper,left=1.45cm,right=1.45cm,top=1.55cm,bottom=1.65cm]{geometry}
\usepackage{amsmath,amssymb,bm}
\usepackage{graphicx}
\usepackage{booktabs,tabularx,array,longtable}
\usepackage{enumitem}
\usepackage[dvipsnames,svgnames]{xcolor}
\usepackage{tikz}
\usetikzlibrary{arrows.meta,positioning,calc,fit,shapes.geometric}
\usepackage[most]{tcolorbox}
\usepackage{hyperref}
\usepackage{caption}
\usepackage{float}
\usepackage{fancyhdr}
\usepackage{titlesec}
\usepackage{microtype}
\IfFontExistsTF{Times New Roman}{\setmainfont{Times New Roman}}{\setmainfont{Liberation Serif}}
\IfFontExistsTF{Noto Sans CJK SC}{\setCJKmainfont{Noto Sans CJK SC}}{}
\IfFontExistsTF{Noto Sans CJK SC}{\setCJKsansfont{Noto Sans CJK SC}}{}
\IfFontExistsTF{DejaVu Sans Mono}{\setmonofont{DejaVu Sans Mono}}{\setmonofont{Latin Modern Mono}}
\definecolor{mainblue}{HTML}{0B5CAD}
\definecolor{softblue}{HTML}{EAF4FF}
\definecolor{darktext}{HTML}{1F2937}
\definecolor{softgray}{HTML}{F5F7FA}
\definecolor{accent}{HTML}{F59E0B}
\definecolor{goodgreen}{HTML}{2E7D32}
\definecolor{badred}{HTML}{B91C1C}
\definecolor{purple}{HTML}{6D5BD0}
\hypersetup{colorlinks=true,linkcolor=mainblue,urlcolor=DarkOrange,citecolor=mainblue}
\setlength{\parindent}{2em}
\setlength{\headheight}{14pt}
\setlength{\parskip}{0.22em}
\setlist[itemize]{leftmargin=2em,itemsep=0.23em,topsep=0.25em}
\setlist[enumerate]{leftmargin=2.2em,itemsep=0.23em,topsep=0.25em}
\renewcommand{\arraystretch}{1.2}
\captionsetup{font=small,labelfont=bf}
\titleformat{\section}{\Large\bfseries\sffamily\color{mainblue}}{\thesection}{0.7em}{}
\titleformat{\subsection}{\large\bfseries\sffamily\color{darktext}}{\thesubsection}{0.6em}{}
\titleformat{\subsubsection}{\normalsize\bfseries\sffamily\color{darktext}}{\thesubsubsection}{0.55em}{}
\pagestyle{fancy}
\fancyhf{}
\fancyhead[L]{\small\sffamily Chain-of-Thought Is Not Always Faithful}
\fancyhead[R]{\small\sffamily \leftmark}
\fancyfoot[C]{\small\thepage}
\renewcommand{\headrulewidth}{0.4pt}
\newtcolorbox{infobox}[1]{enhanced,breakable,colback=softblue,colframe=mainblue!70!black,boxrule=0.8pt,arc=2mm,left=1.5mm,right=1.5mm,top=1mm,bottom=1mm,title=\bfseries #1}
\newtcolorbox{keybox}[1]{enhanced,breakable,colback=softgray,colframe=darktext!45,boxrule=0.6pt,arc=2mm,left=1.5mm,right=1.5mm,top=1mm,bottom=1mm,title=\bfseries #1}
\newtcolorbox{notebox}[1]{enhanced,breakable,colback=accent!8,colframe=accent!80!black,boxrule=0.7pt,arc=2mm,left=1.5mm,right=1.5mm,top=1mm,bottom=1mm,title=\bfseries #1}
\newcommand{\term}[1]{\textbf{\textcolor{mainblue}{#1}}}
\begin{document}

\section{论文基本信息}
\begin{center}
\begin{tabularx}{0.94\linewidth}{>{\bfseries}p{0.18\linewidth}X}
\toprule
题目 & Chain-of-Thought Is Not Always Faithful\\
课件日期 & 2025-07-08\\
研究方向 & 思维链 / 可解释性 / 推理安全\\
一句话概括 & 分析语言模型生成的思维链为何可能只是事后合理化,而非真实内部计算过程,并讨论如何评估与提升其忠实度。\\
\bottomrule
\end{tabularx}
\end{center}

\begin{infobox}{摘要要点}
分析语言模型生成的思维链为何可能只是事后合理化,而非真实内部计算过程,并讨论如何评估与提升其忠实度。 本说明按“研究问题—核心机制—基础公式—实验阅读—局限与优化”的顺序重新组织，不保留原始材料的封面、目录和页码对应。
\end{infobox}

\section{研究问题与动机}
这项工作的出发点不是简单地替换某个网络层，而是识别现有范式中最影响\term{质量、效率、可靠性或可部署性}的瓶颈。结合论文所讨论的问题，可以将研究动机概括为以下几点：
\begin{itemize}
\item 如何平衡 CoT 的有用性和局限性
\item CoT 作为可解释性技术的现状
\item 提升 CoT 保真度的研究方向
\item 实证研究:CoT 的不可靠性
\item CoT 与内部计算的偏差原因
\item CoT 的定义:语言模型在给出最终答案之前生成的多步推理过程。
\end{itemize}

\begin{notebox}{理解重点}
阅读这类论文时，应把“问题是否真实存在”“方法是否直接解决该问题”“实验是否排除了其他解释”分开判断。仅凭更大的模型、更多数据或更长训练得到的提升，不能自动视为结构创新带来的收益。
\end{notebox}

\section{核心思想与整体流程}
分析语言模型生成的思维链为何可能只是事后合理化,而非真实内部计算过程,并讨论如何评估与提升其忠实度。 从信息流角度看，其基本流程可以整理为下图。

\begin{figure}[H]
\centering
\begin{tikzpicture}[>=Latex,node distance=0.78cm]
\tikzstyle{procstep}=[draw,rounded corners,fill=softblue,align=center,minimum width=3.05cm,minimum height=0.95cm,text width=2.75cm]
\node[procstep] (a) {问题输入};
\node[procstep,right=of a] (b) {内部隐式计算};
\node[procstep,right=of b] (c) {显式思维链};
\node[procstep,right=of c] (d) {答案与忠实度评估};
\draw[->,thick,mainblue] (a)--(b);
\draw[->,thick,mainblue] (b)--(c);
\draw[->,thick,mainblue] (c)--(d);
\end{tikzpicture}
\caption{核心信息流与处理阶段。}
\end{figure}

\subsection{方法组成与信息流}
\begin{itemize}
\item 分布式计算与序列化推理的矛盾 ： Transformer 架构的分布式计算与 CoT 的序列化推理之间的矛盾。
\item 作者的观点: ： CoT 解释可能给人以透明度的假象,尤其是在高风险领域,这可能导致对模型输出的 过度信任。
\item ： 应采用更严格的方法来测试解释的保真度,并开发结合 CoT 与因果验证的新方法,以 提高解释的可靠性。
\item 增强人类对 AI 推理的监督 ： 开发更好的工具和框架,帮助人类评估和监督 AI 推理过程。
\item 作者的贡献: ： 综合不同实证研究结果,证明 CoT 不忠实现象是系统性的,而非偶然。
\item ： 探讨 CoT 解释与内部计算偏差的原因,包括分布式处理和人类理性化行为的类比。
\item 未来的研究需要结合 CoT 与更深入的机制解释,以提高模型的可解释性和可 靠性。
\end{itemize}

\begin{keybox}{结构解析}
\begin{itemize}
\item 13 CoT 与内部计算的偏差原因 ； 分布式计算与序列化推理的矛盾 ： Transformer 架构的分布式计算与 CoT 的序列化推理之间的矛盾
\item ； 模型内部存在多个冗余的计算路径 ： 模型可以通过不同的内部路径得出相同的结论
\item ； 生成忠实解释的挑战 ： 即使通过训练干预，模型仍然可能生成不忠实的 CoT
\end{itemize}
\end{keybox}

下图补充展示模块连接、数据流或训练阶段之间的关系。
\begin{figure}[H]
\centering
\includegraphics[width=0.94\linewidth,height=0.43\textheight,keepaspectratio]{method_diagram.jpg}
\caption{模型结构、数据流或主要训练阶段。}
\end{figure}

\section{数学与机制理解}
本节给出理解该工作的基础数学结构。公式用于解释方法所属范式，并不替代原论文中的完整推导。
\begin{equation}
p(y_{1:T}\mid x)=\prod_{t=1}^{T}p(y_t\mid x,y_{<t})
\end{equation}
自回归推理逐 token 生成中间步骤和答案。并行推理、思维链评估与视觉规则归纳都在重新设计这一分解方式，或检验显式文本是否忠实反映内部计算。

\begin{keybox}{从公式回到方法}
应重点观察论文究竟改变了公式中的哪一部分：是输入表示、连接矩阵、条件变量、参数化方式、训练目标，还是求解与调度过程。真正的创新通常可以在这一层被准确定位。
\end{keybox}


\section{实验结果应如何阅读}
\begin{itemize}
\item 尽管 CoT 被广泛使用,但其可解释性并未得到充分验证。
\item 音频处理:用于音频情感识别、说话人意图分类等。
\item CoT 在多个领域被广泛用于提升可解释性。
\item 视觉任务:用于自动驾驶、情感识别等领域。
\item 医疗诊断:用于解释医学图像诊断结果。
\item 法律应用:用于法律推理和判决预测。
\item CoT 作为可解释性技术的现状
\end{itemize}

\begin{notebox}{结果解读}
\begin{itemize}
\item 9 作者的贡献与观点 ； 作者的贡献： ： 综合不同实证研究结果，证明 CoT 不忠实现象是系统性的，而非偶然
\item ： 探讨 CoT 解释与内部计算偏差的原因，包括分布式处理和人类理性化行为的类比
\item ： 识别特定条件（如提示偏差、复杂多步推理和预设答案）下不忠实现象最为普遍
\item ； 作者的观点： ： CoT 解释可能给人以透明度的假象，尤其是在高风险领域，这可能导致对模型输出的 过度信任
\end{itemize}
\end{notebox}

该部分以文字概括实验结论与比较条件，避免用整页截图代替分析。
实验部分不应只看单一最好数字。还需要核对模型规模、训练数据、训练步数、采样或解码预算、硬件条件以及是否使用额外蒸馏、检索或后处理。只有比较条件接近时，结果才能真正说明方法本身的优势。

\section{优点、局限与可优化空间}
\subsection{主要优点}
\begin{itemize}
\item \term{问题与方法对应明确}：论文围绕一个可识别的核心瓶颈展开，方法结构能够直接映射到该瓶颈。
\item \term{可复用的设计思想}：即使不完整复现模型，其中的分解、稀疏化、递归、条件控制或系统优化思想也可迁移到相邻任务。
\item \term{实验提供了多角度证据}：实验通常同时展示主结果、效率比较、案例或消融，使结论不完全依赖单一指标。
\end{itemize}

\subsection{局限与进一步优化}
\begin{itemize}
\item 需要进一步区分\term{结构收益}与更大算力、更多数据、不同训练技巧带来的收益，并在等参数、等计算预算下进行严格比较。
\item 对超参数、输入分布和模型规模的敏感性仍应系统分析；在一个数据集上的最优配置不一定能直接迁移。
\item 实际部署还要考虑显存峰值、并发、延迟抖动、失败案例和鲁棒性，而不仅是离线平均指标。
\item 可尝试将该方法与蒸馏、量化、缓存复用、动态路由或更强的表示学习结合，验证不同优化是否互补。
\end{itemize}

\section{结论}
总体而言，分析语言模型生成的思维链为何可能只是事后合理化，而非真实内部计算过程，并讨论如何评估与提升其忠实度。。 进一步需要注意：如何平衡 CoT 的有用性和局限性;CoT 作为可解释性技术的现状。

\end{document}
